%%%%%%%%%%  scan idx 14  |  folio 1  %%%%%%%%%%
\section*{\underline{Exposé 0}}
\begin{center}
\underline{\underline{ESQUISSE D'UN PROGRAMME POUR UNE THÉORIE DES INTERSECTIONS}}\\[3pt]
\underline{\underline{SUR LES SCHÉMAS GÉNÉRAUX}}\\[8pt]
\underline{par A. Grothendieck}
\end{center}

Le présent exposé est de nature introductif, et sa lecture n'est pas logiquement indispensable pour l'étude du Séminaire. Il s'adresse plus particulièrement aux lecteurs au courant de la version provisoire du théorème de Riemann--Roch, telle qu'elle est exposée dans le rapport de Borel--Serre [2] ou dans le rapport de Grothendieck déjà mentionné dans l'Introduction (cité [RRR] par la suite), et qui est reproduit sous forme d'appendice à la fin du présent exposé.

\subsection*{1.}
Rappelons la formule de Riemann--Roch pour un morphisme propre
\[
f:X.\longrightarrow Y
\]
de schémas quasi-projectifs et lisses sur un corps $k$, et un faisceau cohérent $F$ sur $X$, dont la classe dans le groupe $K(X)$ des classes de faisceaux cohérents sur $X$ est notée par $\cl(F)$~:
\begin{equation}\tag{1.1}
\Todd(T_Y)\,\ch_Y\bigl(f_*(\cl(F))\bigr)=f_*\bigl(\Todd(T_X)\,\ch_X(F)\bigr).
\end{equation}
Cette formule est valable dans $A(Y)\otimes_{\mathbb Z}\mathbb Q$, où $A(Y)$ est l'anneau de Chow de $Y$, le $f_*$ du second membre étant déduit par tensorisation par $\mathbb Q$ de l'homomorphisme «~image directe de cycles~»
\[
f_*:A(X)\longrightarrow A(Y),
\]
celui du premier étant la caractéristique d'Euler--Poincaré de $F$ relativement à $f$~:
\[
f_*(\cl(F))=\sum_i(-1)^i\cl\bigl(R^if_*(F)\bigr)~;
\]
$\ch_X$, $\ch_Y$ désignent les caractères de Chern sur $X$ resp. sur $Y$, et $T_X$, $T_Y$ les fibrés tangents à $X$ resp. à $Y$. Comme on sait, $\Todd(-)$ et $\ch(-)$ sont des polynômes universels à coéfficients dans $\mathbb Q$ en les classes de Chern de l'argument. Le terme constant de $\Todd(-)$ étant $1$, c'est un élément inversible pour toute valeur de l'argument, de sorte que la formule (1.1) prend après multiplication par $\Todd(T_Y)^{-1}$ la forme équivalente, plus commode pour notre propos
\begin{equation}\tag{1.2}
\ch_Y\bigl(f_*(\cl(F))=f_*\bigl(\Todd(T_f)\cdot\ch_X(F)\bigr),
\end{equation}
où on a posé
\begin{equation}\tag{1.3}
T_f=T_X-f^*(T_Y)\in K(X),
\end{equation}
